\section{Applications of the method to generic composite operators}
\label{sec:casivari}

In the previous section, using the renormalization of the scalar and
pseudoscalar densities as a prototype,
the general strategy of the non-perturbative scheme has been presented.
 In the present section the necessary modifications
needed in other cases are briefly described. We will follow the
classification of operators   introduced in sec.\ \ref{sec:idea}.

\subsection{Finite operators}

In the case of vector and axial vector
currents, we do not have the freedom of  fixing the
renormalization conditions in an arbitrary way:
the renormalization conditions are acceptable only
when  they are compatible with the Ward identities.
To be specific, one can apply  eq.\ (\ref{eq:zdef}) to
 the local vector current in order
to determine $Z_{V^L}$
\beq Z_{V^L} = \frac{Z_\psi}{\Gamma_{V^L}}, \label{eq:zvl} \eeq
We now show  that this definition of $Z_V^L$,
which uses  the  general prescription given
in the previous section,  coincides
with the one derived from the vector current Ward identity  on quark
states \cite{bo,ka}\footnote{For the remainder of this subsection we work
in lattice units, setting $a=1$.}
\beq
q_{\mu} \Bigl[Z_{V^L} \Lambda_{V^L_\mu}(p+\frac{q}{2},p-\frac{q}{2})
\Bigr]=-i\Bigl( S^{-1}(p+\frac{q}{2}) - S^{-1}(p-\frac{q}{2})\Bigr).
\label{eq:pallino}\eeq
where $\Lambda_{V^L_\mu}(p+q/2,p-q/2)$ is the amputated
vertex with momentum transfer $q$.
By applying $\partial / \partial q_\rho$, at $q=0$, on both sides of eq.
(\ref{eq:pallino}), one gets
\beq
Z_{V^L}\Bigl( \Lambda_{V^L_\rho}(p) + q_\mu \frac{\partial}{\partial q_\rho}
 \Lambda_{V^L_\mu}(p+\frac{q}{2},p-\frac{q}{2})\vert_{q=0}\Bigr)
=-i \frac{\partial}{\partial p_\rho}S^{-1}(p).
\label{eq:pinco} \eeq
The second term on the l.h.s.\ vanishes when $q \to 0$.
By tracing eq.\ (\ref{eq:pinco}) with $\gamma_{\rho}$,
we then recover the definition of $Z_{V^L}$ of eq.\ (\ref{eq:zvl}).

One is tempted to follow the same path for the local axial vector current,
i.e.\ to define
\beq Z_{A^L} = \frac{Z_\psi}{\Gamma_{A^L}}.  \label{eq:zal} \eeq
However this definition of $Z_{A^L}$   is not equivalent to the one
obtained by the Ward identity
\beq
q_{\mu} \Bigl[Z_{A^L} \Lambda_{A^L_\mu}(p+\frac{q}{2},p-\frac{q}{2})
\Bigr]=-i\Bigl( \gamma_5 S^{-1}(p+\frac{q}{2}) + S^{-1}(p-\frac{q}{2})
 \gamma_5 \Bigr).
\label{eq:awi}\eeq
To see this we proceed as for the vector current,  applying $\partial /\partial
 q_\rho$
to both sides of eq.\  (\ref{eq:awi}).
\beq Z_{A^L}\Bigl( \Lambda_{A^L_\rho}(p) + q_\mu
\frac{\partial}{\partial q_\rho}
 \Lambda_{A^L_\mu}(p+\frac{q}{2},p-\frac{q}{2})\vert_{q=0}\Bigr) =
  \nn \eeq
\beq  -\frac{i}{2}\Bigl( \gamma_5 \frac{\partial}{\partial p_\rho}S^{-1}(p)
-\frac{\partial}{\partial p_\rho}S^{-1}(p) \gamma_5 \Bigr).
\label{eq:dawi} \eeq
 In the axial case, however,  the second term
 on the l.h.s.\ does not vanish, as $q \to 0$ and, indeed,
this term is necessary to saturate
the Ward identity, in presence of a massless Goldstone boson.
As demonstrated in the appendix however, this term becomes negligible
for high values of $p^2$. In this case, by tracing eq.\ (\ref{eq:dawi}),
one recovers (\ref{eq:zal}). In conclusion,
we expect to find a value of
$Z_{A^L}$ close to those obtained from the Ward identity on hadron or
 on quark states \cite{wi}, by using eq.\ (\ref{eq:zal}), at large
$\mu^2$.

$Z_S/Z_P$ is another finite quantity, which is completely fixed by
the Ward identities, even though the renormalization conditions
on $Z_S$ and $Z_P$  are separately arbitrary. In other words
one is free to  decide the renormalization conditions
of one of the two operators,
the scalar density  say, and the value of $Z_P$ will automatically follow.
The procedure of sec.\ \ref{sec:idea}, respects the Ward identities
and thus  guarantees that the correct value of $Z_S/Z_P$ is obtained.
This is clearly true only at large values of $\mu^2$,
for the same reasons as for
the axial vector current.

\subsection{Logarithmically divergent operators}

For this category of operator only the question of operator-mixing is left to
be discussed.
There are two kinds of mixing: the one which can occur also in
the continuum and the one induced by the explicit lattice
chiral symmetry breaking.
We are only  concerned  here with the latter case,
and hence restrict our discussion to operators
which renormalize multiplicatively in the continuum.
To illustrate our arguments, we consider  the four-fermion operator
\beq O^{\Delta S=2} = \Bigl(\bar s \gamma^\mu (1-\gamma_5) d \Bigr)
\Bigl(\bar s \gamma_\mu (1-\gamma_5) d \Bigr), \label{eq:ds2b} \eeq
which is relevant to the calculation of the $K^0$--$\bar K^0$ mixing
amplitude.
Because of the Wilson term, this operator mixes with operators of different
chiralities \cite{bo,te,4ferm}
\beq  O^{\Delta S=2}_{{\rm cont}} = Z_O \Bigl( O^{\Delta S=2}_{{\rm
latt}} + \sum_i Z_i O^i \Bigr),\label{eq:omix} \eeq
where the operators $O^i$ are given by $O^1= (\bar s d )(\bar s d),
O^2=(\bar s
\gamma_5 d)(\bar s \gamma_5 d)$, etc.\ \cite{4ferm}.
 The mixing coefficients $Z_i$ are finite
functions of $g_0^2(a)$.
The over-all renormalization $Z_O$ is necessary to eliminate
the logarithmic divergence.

We now sketch the renormalization procedure for $O^{\Delta S=2}$.
Let us start from the four-point amputated, Green function of
$O^{\Delta S=2}_{{\rm latt}}$, between quark states with momentum $p^2=\mu^2$
\beq \Lambda_{\alpha \beta \gamma \delta}^{ABCD}(pa),
\label{eq:g4} \eeq
where $\alpha, \beta \dots$ and $A, B \dots$ are spin and colour
indices respectively.
To find the mixing coefficients $Z_i$, one uses suitable projectors
$\hat P$, as
was done for two-quark operators\footnote{ In order to simplify the
presentation,  we do not discuss here the
complications arising from the different possible colour structures.},
for example
\beq {\rm tr } ( \Lambda \hat P )= \Lambda_{ \alpha \beta
 \gamma \delta }^{AABB}( \gamma_5)_{ \beta \alpha}
 ( \gamma_5 )_{ \delta \gamma}. \eeq
By projecting on all the possible colour-spin structures, corresponding
to the operators $O^i$, one obtains a system of linear equations in
the mixing coefficients $Z_i$.  Using the $Z_i$, determined in this way,
one can compute the Green function of the subtracted operator
$O^s=O^{\Delta S=2}_{{\rm latt}}+ \sum_i Z_i O^i$. From the amputated vertex of
$O^s$, one can finally remove the logarithmic divergence, by imposing the
renormalization condition
\beq Z_O \Bigl( \mu , g(a) \Bigr) Z^{-2}_{ \psi }
 \Bigl( \mu , g(a) \Bigr) \Gamma^s_O ( pa) \vert_{p^2 =
 \mu^2} = 1, \eeq
where
\beq \Gamma^s_O \sim { \rm tr } (
 \Lambda^s \hat P_{ L \times L} ) ,\eeq
and $P_{ L \times L}$ is the suitable projector on the
$ \gamma_ \mu ( 1 - \gamma_5 ) \otimes \gamma^\mu (
1 - \gamma_5 ) $ Dirac tensor.

\subsection{Power divergent operators}

Composite operators are divergent in $1/a$  when they can mix with lower
dimensional operators. Examples are given by the
operators relevant to deep inelastic scattering and by
the penguin-operators  appearing in the $\Delta S=1/2$ Hamiltonian.
Let us consider the $\Delta S=1/2$ operator
\beq O^{ \Delta S=1/2}= ( \bar s \gamma_\mu ( 1 - \gamma_5 )
d ) \sum_i (\bar q_i \gamma^\mu ( 1 + \gamma_5 ) q_i) .\label{eq:dt12} \eeq
$ O^{ \Delta S=1/2}_{{\rm latt}}$ mixes with the scalar and pseudoscalar
denisties, even at zeroth order in $\alpha_s$, through the diagram in
fig.\ \ref{fig:penguin}\footnote{ $O^{\Delta S=1/2}_{{\rm latt}}$ also  mixes
 with the chromo-magnetic operator, e.g.\ $\bar s \sigma_{\mu \nu}
G^{\mu \nu} d$, which is not being considered here.}
\beq O^{\Delta S=1/2}_{{\rm cont}} =
Z_O O^{\Delta S=1/2}_{{\rm latt}}+ Z_1 (\bar s d )+ Z_2 (\bar s  \gamma_5
 d )+ \dots. \label{eq:ds12} \eeq
\begin{figure}[htbp]
\centering
\includegraphics[width=0.62\textwidth]{images/feynman1.pdf}
\caption{Diagram which mixes the four fermion operator in
eq.\ (\ref{eq:dt12}) with the scalar
and pseudoscalar densities. The mixing is induced by the chiral
breaking term present in the Wilson and Clover actions.}
\label{fig:penguin}
\end{figure}

To find the coefficients $Z_{1,2}$, one imposes that the two-point
Green function of the operator on the r.h.s.\ of eq.
(\ref{eq:ds12}), on off-shell quark states,  is zero \cite{ber1}.
Once  all the divergent parts have been removed, one can impose the
renormalization conditions, necessary to remove the over-all logarithmic
divergence, as discussed above.
